\documentclass[11pt]{article}
\pdfinfoomitdate=1
\pdftrailerid{}
\pdfsuppressptexinfo=15
\usepackage[a4paper,margin=1in]{geometry}
\usepackage[T1]{fontenc}
\usepackage{lmodern}
\usepackage{microtype}
\usepackage{amsmath,amssymb,amsthm,mathtools}
\usepackage{booktabs,array,longtable,enumitem,xcolor}
\usepackage[colorlinks=true,linkcolor=blue!55!black,urlcolor=blue!55!black]{hyperref}

\newtheorem{theorem}{Theorem}[section]
\newtheorem{proposition}[theorem]{Proposition}
\theoremstyle{definition}
\newtheorem{definition}[theorem]{Definition}
\theoremstyle{remark}
\newtheorem{remark}[theorem]{Remark}
\newcommand{\RH}{RH}
\newcommand{\dd}{\mathrm d}
\newcommand{\Eoff}{E_{\rm off}}

\title{\textbf{Conditional Prime Dynamics: a Corpus Frontier Synthesis}\\
\large Provenance, branch separation, and the unpaid theorem budget in RH-1--RH-361}
\author{Liang Wang}
\date{August 2026}

\begin{document}
\maketitle

\begin{abstract}
The conditional prime-dynamics program has produced 361 numbered RH papers,
many of them intentionally narrow reviews, certificates, or scoped negative
results.  This document is a provenance-preserving synthesis rather than a
replacement corpus: the numbered papers remain the atomic source of every
claim.  We compress the corpus into five functional phases, separate the
actual noisy branch from the deterministic counterloop branch, and state the
exact common-clock identities that prevent promotion between them.  The
latest source-locked frontier is the unnormalised same-clock head budget
\[
 D_{4k}(R)=\sum_{2\le n<4k}|h_{\sigma,n}-s_{k,n}|R^n/n\longrightarrow0,
\]
which is assumed by the RH-355--RH-360 transfer statements and proved by
none of the sources.  The synthesis therefore records
\texttt{NOT\_TESTABLE} for the
next numbered route.  All five Gates A--E remain false/open.  No
Hilbert--Polya operator, Riemann-zero identification, von Mangoldt trace,
completed-zeta divisor identity, or proof of the Riemann Hypothesis is
claimed.
\end{abstract}

\section{Purpose and status vocabulary}

This paper changes altitude, not the mathematics.  It is intended to make a
large research record readable before a new idea is attempted.  It does not
delete, rewrite, or silently merge any numbered paper.  A machine-readable
inventory records all unique numbers $1,\ldots,361$ and the four historical
duplicate directory groups.  The inventory hashes selected source files and
checks presence; it is not a re-proof of the corpus.

We use five non-interchangeable labels:
\begin{description}[leftmargin=2.8cm,style=nextline]
\item[proved] an analytic theorem in the scope written in its source paper;
\item[certified] a finite exact or outward-rounded computation;
\item[conditional] a valid implication whose physical hypotheses remain open;
\item[scoped negative] a proved non-implication or obstruction for the named
 route, not for every possible model;
\item[open/NOT\_TESTABLE] the required source, identification, clock, or
 uniform estimate is absent.
\end{description}

\section{Provenance topology}

The filesystem contains 361 unique numerical labels and 365 directories
matching the numerical naming pattern.  The difference is explained by four
empty historical aliases at RH-302, RH-303, RH-304, and RH-306.  The audit
therefore identifies a source by its full canonical directory path and a
SHA-256 digest, never by its number alone.  It chooses the candidate containing
the README, TeX source, and PDF, and records the empty alternative as an alias
rather than deleting it.

There are 29 established review anchors:
\begin{quote}\small
RH-71, 81, 91, 100, 119, 129, 139, 149, 159, 171, 181, 191, 201,
211, 221, 231, 241, 251, 261, 271, 281, 291, 301, 311, 321, 331,
341, 351, and 361.
\end{quote}
Their declared input ranges cover 349 of the 361 numbered IDs.  The twelve
exceptions are RH-101--RH-109, RH-150, RH-160, and RH-161.  They are not
discarded: RH-MVP1 covers RH-1--RH-160, while RH-161 is the independent
packet-to-Riesz relative determinant assembly feeding the next phase.  This
matters because the review batches are not all identical ten-paper windows.
For example, RH-159 audits RH-151--RH-158 and leaves RH-150 separate, while
RH-171 audits RH-162--RH-170 and leaves RH-161 separate.

\begin{proposition}[Provenance-preserving compression]
The canonical inventory is a surjection from 361 numerical labels to 361
nonempty source directories, together with four recorded empty aliases.  The
synthesis changes neither a source hash nor the logical status attached to a
source claim.
\end{proposition}
\begin{proof}
The executable inventory enumerates every numerical label from 1 through 361,
requires a README, TeX source, and PDF in the chosen directory, and hashes the
available README, manuscript, theorem ledger, roadmap, and result record.  The
archive verifier replays all 1,356 selected source hashes.  This proves
presence and identity only; it deliberately does not claim to re-prove the
source theorems.
\end{proof}

\section{Five functional phases}

The numbered sequence can be read as the following dependency map.  The
intervals are navigation ranges, not claims that every paper in an interval
has the same logical strength.

\begin{longtable}{>{\raggedright\arraybackslash}p{1.4cm}>{\raggedright\arraybackslash}p{4.0cm}>{\raggedright\arraybackslash}p{8.0cm}}
\toprule
range & role & durable content and boundary \\
\midrule
\endhead
\RH-1--160 & symbolic, determinant, packet, reset, and MVP architecture &
 corrected sieve--kneading coordinates, parity and determinant hygiene,
 finite/continuum certificates, and a conditional five-interface roadmap;
 no all-level intrinsic determinant, self-adjoint generator, arithmetic trace,
 or RH conclusion. \\
\RH-161--241 & physical Riesz and trace-envelope frontier &
 packet/transport and regularised determinant analyses, culminating in the
 moving noisy all-order envelope question; RH-241 explicitly leaves the
 uniform envelope and coefficient anchor open. \\
\RH-242--271 & deterministic numerator and analytic-tail tools &
 deterministic coefficient anchors, all-order envelopes, sharp radii, and
 quotient/tail certificates; these are target-side results and do not identify
 a moving noisy head. \\
\RH-272--341 & prefix, annular, endpoint, spectral, and alias routes &
 exact typed decompositions, rate-free or conditional diagonal bridges,
 endpoint and alias obstructions, and information-class negatives; no common
 physical full-prefix budget is closed. \\
\RH-342--361 & signed completion and counterloop frontier &
 selected/normalised actual $p,Y$ results, an unconditional deterministic
 counterloop $s$ branch, and the exact typed separation theorem; the physical
 head-defect bridge remains open. \\
\bottomrule
\end{longtable}

\section{What the intermediate frontiers establish}

\subsection{Foundation and Stage-A architecture: RH-1--RH-160}

RH-MVP1 already compresses this range into a rigorous foundation $F$ and five
missing interfaces A--E.  The foundation includes the corrected symbolic
coordinate, parity-resolved deterministic trace geometry, fixed-noise
regularised determinants, continuum bridges, and finite certification tools.
Its conditional spectral-divisor closure is a valid implication, but none of
its five macro premises is supplied merely by the 160-layer audit.  The first
unpaid interface is still the canonical all-level determinant A.

\subsection{Physical clouds and trace envelopes: RH-161--RH-241}

This range passes through packet-to-Riesz assembly, history/cycle identities,
biorthogonal clocks, source-channel quartets, divisor-first transport, and
rank-growing clouds.  The later batches establish finite and local structure
for projection-free regularised determinants.  RH-241 records 7,280 finite
ledger items and a subunit trace envelope for orders 2--12, but explicitly
leaves the moving all-order noisy envelope and coefficient identification
open.  Repeating or extending the finite atlas is not an all-order theorem.

\subsection{Deterministic target closure: RH-242--RH-281}

The superloop and selector papers isolate which finite anchored completions
are unavailable.  RH-263 gives the deterministic numerator coefficient
anchor; RH-267 gives a certified unified all-order deterministic trace
envelope; RH-268 gives the sharp deterministic coefficient radius.  RH-281
adds the deterministic counterloop bridge.  These are genuine all-order
target-side results, but they do not identify the moving noisy cloud, actual
head, or full-trace coefficient.

\subsection{Noisy tails, annuli, and endpoint spectra: RH-282--RH-321}

This range contains genuine finite modulus-complete noisy heads, analytic
complement tails, weighted-prefix clock thresholds, and exact annular
criteria.  It also gives sharp synthetic spectral endpoint realisability and
saturation results.  The distinctions are decisive: RH-294 has a rate-free
weighted full-trace diagonal on an unspecified slow clock; RH-300 is a
sufficient annular criterion; and RH-321 is sharp for the declared synthetic
spectral class.  None identifies the physical logarithmic clock with the
complete actual head required at the later first-alias frontier.

\subsection{First alias and signed completion: RH-322--RH-361}

RH-322--RH-341 build exact local physical formulas and transfer criteria but
leave the aggregate actual obligations open.  RH-342--RH-351 then prove
information-class non-identifiability and signed-completion results; these
are scoped negatives, not pairs of physical noisy operators.  Finally,
RH-352--RH-361 isolate the actual selected/normalised branch from the
unconditional deterministic counterloop branch.  This is the current endpoint
used below.

\section{The current two-branch statement}

On the source-locked Hardy clock
\[
 k=\frac{\log(1/\sigma)}{2\log\lambda}+O(1),\qquad 2\le n<4k,
\]
write $p$ for the actual direct coefficient, $q$ for the full-trace
coefficient, $h$ for the actual modulus-complete head, and $s$ for the
deterministic graded counterloop.  The source papers give the typed identities
\begin{equation}
 p=\tau-a=q-d,\qquad d=h-s,
 \qquad q=p+d,\qquad h=s+d. \label{eq:typed}
\end{equation}

The actual branch (RH-352--RH-354) controls selected or normalised $p$ and
$Y$ quantities.  The deterministic branch (RH-355--RH-360) closes a sequence
of budgets for $s$, including terminal localisation, logarithmic inversion,
and an exponential-tilt phase diagram.  Equation~\eqref{eq:typed} is an
identity, not an inclusion of spectra and not a physical realisation of $s$.

\begin{proposition}[Typed branch separation]
The information supplied by the current corpus does not imply a physical
bound for $q$ or $h$ from the proved selected/normalised $p$ results and the
deterministic $s$ results alone.  Any such promotion requires an independent
estimate of the defect $d$ on the same clock and order range.
\end{proposition}
\begin{proof}
The identities in \eqref{eq:typed} show that both promotions contain the
unpaid term $d$.  RH-361 gives the finite coefficient-fibre calculation
\(d[e]=e\), \(q[e]=p+e\), and \(h[e]=s+e\) for arbitrary signed arrays $e$.
This is an information-class non-implication.  The construction does not
assert that every fibre is realised by a noisy operator; therefore it is not a
physical counterexample and cannot be used to promote or refute a spectral
claim.
\end{proof}

\section{The first missing leaf}

Every actual-head inheritance statement in RH-355--RH-360 assumes the
unnormalised budget
\begin{equation}
 D_{4k}(R)=\sum_{2\le n<4k}|h_{\sigma,n}-s_{k,n}|R^n/n\to0. \label{eq:defect}
\end{equation}
The selected windows and normalisations in RH-352--RH-354 do not pay the
loss in \eqref{eq:defect}; absolute majorants do not replace signed aggregate
cancellation.  RH-287 and RH-294 provide slower or unspecified diagonal
results without identifying the physical clock.  RH-241 still lacks the
moving noisy all-order envelope and its no-over-extraction coefficient bridge.

\begin{theorem}[Repository frontier theorem]
Relative to the source files and claim firewall recorded in
\texttt{RH\_HANDOFF.md} at the current endpoint RH-361, no admissible source
proves \eqref{eq:defect}, proves a genuine physical obstruction to it, or
supplies an independent same-type full-trace $q/\Eoff$ theorem on the common
clock.  Consequently RH-362 is not activated by the existing corpus.
\end{theorem}
\begin{proof}
The RH-361 review and its source-lock/auditor records list RH-355--RH-360 as
conditional on \eqref{eq:defect}; they explicitly distinguish the normalised
upper-band quantity from the unnormalised full prefix.  The alternative-route
scan records that RH-287/294 use a different diagonal clock, while RH-334,
RH-339, RH-344, RH-346, and RH-348 give typed identities without an aggregate
bound.  The remaining admissible triggers in the handoff are exactly these
missing leaves.  This is a repository audit theorem, not a claim that no
mathematical route exists outside the corpus.
\end{proof}

\section{Durable non-promotions}

Several recurring shortcuts are already excluded by source type or by a
proved scoped negative result.  A synthesis must keep them visible:

\begin{center}
\small
\begin{tabular}{>{\raggedright\arraybackslash}p{4.1cm}>{\raggedright\arraybackslash}p{8.7cm}}
\toprule
invalid promotion & unpaid information \\
\midrule
finite rows $\Rightarrow$ all-order theorem & a uniform moving-order estimate and asymptotic quantifiers \\
fixed or synthetic spectrum $\Rightarrow$ actual noisy operator & a physical realisation and identification theorem \\
slow diagonal clock $\Rightarrow$ logarithmic physical clock & a quantitative lower clock bound and common order range \\
normalised tail $\Rightarrow$ unnormalised prefix & the removed exponential scale and all omitted orders \\
selected window $\Rightarrow$ complete $\Eoff$ & odd, lower, critical, and upper-alias coordinates on one type \\
separate absolute bounds $\Rightarrow$ signed aggregate obstruction & a signed cancellation or lower-bound theorem \\
coefficient fibre $\Rightarrow$ physical counterexample & realisability by noisy operators with the required ranks and roots \\
deterministic $s$ $\Rightarrow$ actual $h$ or $q$ & the common-clock defect $d=h-s$ \\
\bottomrule
\end{tabular}
\end{center}

\section{Gate matrix and claim firewall}

\begin{center}
\begin{tabular}{>{\raggedright\arraybackslash}p{1.0cm}>{\raggedright\arraybackslash}p{4.0cm}>{\raggedright\arraybackslash}p{2.0cm}>{\raggedright\arraybackslash}p{6.2cm}}
\toprule
Gate & object & status & reason for non-promotion \\
\midrule
A & canonical intrinsic dynamical determinant & false/open & no typed all-level moving-cloud assembly \\
B & order-sensitive scattering/unitary completion & false/open & no canonical orientation-preserving completion \\
C & self-adjoint generator and intrinsic $T\log T$ law & false/open & no proved generator or counting law \\
D & von Mangoldt-weighted prime-power trace & false/open & no target-independent arithmetic trace identity \\
E & completed-zeta spectral divisor equality & false/open & no no-missing/no-spurious divisor theorem \\
\bottomrule
\end{tabular}
\end{center}

In particular, finite rows are reproduction checks; synthetic spectra are not
physical noisy operators; conditional criteria are inactive until their
hypotheses hold on one common clock; and deterministic shells are not actual
root sets, ranks, or spectral submultisets.  The synthesis does not construct
a Hilbert--Polya operator, identify Riemann zeros, prove the von Mangoldt
trace, prove completed-zeta divisor equality, or prove RH.

\section{What would justify a new numbered paper?}

The shortest admissible reopen triggers are:
\begin{enumerate}[leftmargin=2.2em]
\item a physical same-clock proof of \eqref{eq:defect}, or a genuine physical
 obstruction to it;
\item a typed actual full-trace $q$ or complete $\Eoff$ aggregate theorem;
\item a complete unnormalised direct prefix including orders below the
 RH-354 moving cut;
\item the RH-241 moving noisy all-order envelope plus a valid coefficient
 bridge; or
\item another independent source-backed theorem edge.
\end{enumerate}
More computation of finite rows, reparameterisation of RH-358--RH-360,
normalisation changes, or relabelling a coefficient called $q$ does not meet
one of these triggers.

\section{One umbrella and four long-form volumes}

The present manuscript is the short umbrella: it fixes vocabulary,
provenance, phases, Gates, and the current stop.  The long-form publication
layer consists of four independently auditable volumes without changing any
theorem status:
\begin{enumerate}[leftmargin=2.2em]
\item Volume I, RH-1--RH-160: foundations and the conditional A--E
 architecture, retained in RH-MVP1;
\item Volume II, RH-161--RH-241: physical Riesz packets, temporal clouds, relative
 determinants, and the trace-envelope frontier;
\item Volume III, RH-242--RH-281: deterministic numerator anchors, selectors, analytic
 tails, and counterloops;
\item Volume IV, RH-282--RH-361: noisy heads, weighted/annular endpoints, first alias,
 and actual-versus-deterministic signed completion.
\end{enumerate}
The four ranges are disjoint and cover RH-1--RH-361.  This division is
preferable to one purported mega-theorem because each volume exposes its data
types, clocks, and independently auditable boundary.

\section{Conclusion}

The corpus can therefore be treated as one coherent research map without
pretending that it is one theorem.  The useful compression is semantic:
preserve the individual sources, expose the two typed branches, record the
unpaid defect budget, and stop the numbered sequence until a new edge appears.
That leaves room for genuinely new ideas while preventing old finite or
conditional evidence from being recycled as a full-order result.

\end{document}
